\section*{Capítulo \ref{ch:g12:curly-arrows} --- Mecanismos de reacción: flechas curvas}

\begin{solution}{exo:g12:curly-arrows:1}
\ce{OH-}: dador (\omterm{def:g10:lewis-and-shape:lone-pair}{pares libres}, carga negativa); ningún sitio aceptor.
\ce{H+}: aceptor (no tiene ningún \omterm{def:g9:inside-the-atom:nucleus}{electrón}). \ce{CH2=CH2}: dador (el
\omterm{def:g10:lewis-and-shape:multiple-bond}{doble enlace}). \ce{CH3-Cl}: dador (\omterm{def:g10:lewis-and-shape:lone-pair}{pares libres} del \ce{Cl}), aceptor (el
carbono, $\delta^+$). \ce{H2O}: dador (\omterm{def:g10:lewis-and-shape:lone-pair}{pares libres} del \ce{O}); sus
\omterm{def:g7:atoms-and-molecules:atom}{átomos} de hidrógeno, $\delta^+$, son sitios aceptores débiles.
\end{solution}

\begin{solution}{exo:g12:curly-arrows:2}
(a) \omterm{def:g12:curly-arrows:categories}{sustitución}; (b) \omterm{def:g12:curly-arrows:categories}{adición}; (c) \omterm{def:g12:curly-arrows:categories}{eliminación}; (d) \omterm{def:g12:curly-arrows:categories}{adición}.
\end{solution}

\begin{solution}{exo:g12:curly-arrows:3}
Empieza en un par de \omterm{def:g9:inside-the-atom:nucleus}{electrones} (un \omterm{def:g10:lewis-and-shape:lone-pair}{par libre} o un enlace) y termina en
el \omterm{def:g7:atoms-and-molecules:atom}{átomo} al que va el par. De un \omterm{def:g10:lewis-and-shape:lone-pair}{par libre} del oxígeno a un carbono: el
par pasa a ser un nuevo enlace \ce{O-C}.
\end{solution}

\begin{solution}{exo:g12:curly-arrows:4}
El amoníaco cede el \omterm{def:g10:lewis-and-shape:lone-pair}{par libre} de su nitrógeno; la flecha va de ese \omterm{def:g10:lewis-and-shape:lone-pair}{par
libre} al \ce{H+}. El producto \ce{NH4+} lleva la carga $+1$.
\end{solution}

\begin{solution}{exo:g12:curly-arrows:5}
Una especie formada en una etapa y consumida en otra posterior; no
figura en la ecuación global. Aquí, el carbocatión \ce{CH3-CH2+}.
\end{solution}

\begin{solution}{exo:g12:curly-arrows:6}
Una flecha de un \omterm{def:g10:lewis-and-shape:lone-pair}{par libre} del oxígeno del \ce{OH-} al carbono; otra del
enlace \ce{C-Br} al bromo. Cargas: $-1$ a la izquierda (\ce{OH-}); a la
derecha, el metanol es neutro y el \ce{Br-} lleva $-1$: el total es $-1$
en ambos lados.
\end{solution}

\begin{solution}{exo:g12:curly-arrows:7}
El bromo es más electronegativo que el hidrógeno: se queda con el par y
sale como \ce{Br-} ($-1$); el hidrógeno se une a un carbono, y el otro
carbono se queda con una carga positiva: \ce{CH3-CH2+} ($+1$).
\end{solution}

\begin{solution}{exo:g12:curly-arrows:8}
1.~\ce{CH3-CH2-OH + H+ -> CH3-CH2-OH2+}: dador, un \omterm{def:g10:lewis-and-shape:lone-pair}{par libre} del
oxígeno; aceptor, el \ce{H+}. 2.~\ce{CH3-CH2-OH2+ -> CH3-CH2+ + H2O}: el
par \ce{C-O} va al oxígeno positivo (aceptor), que sale con él en forma
de agua. 3.~\ce{CH3-CH2+ -> CH2=CH2 + H+}: el par de un enlace \ce{C-H}
(dador) se desplaza hacia el carbono positivo (aceptor) para formar el
\omterm{def:g10:lewis-and-shape:multiple-bond}{doble enlace}, y el \ce{H+} sale. El \omterm{def:g9:ions:ion}{ion} hidrógeno usado en la etapa 1 se
recupera en la etapa 3: aparece en ambos lados y se simplifica.
\end{solution}

\begin{solution}{exo:g12:curly-arrows:9}
Una flecha de un \omterm{def:g10:lewis-and-shape:lone-pair}{par libre} del oxígeno del agua al carbono positivo. Tras
la pérdida de \ce{H+}, etanol \ce{CH3-CH2-OH}. Global:
\ce{CH2=CH2 + H2O -> CH3-CH2-OH}, una \omterm{def:g12:curly-arrows:categories}{adición}.
\end{solution}

\begin{solution}{exo:g12:curly-arrows:10}
Se desplazan dos pares, en una sola etapa: se forma el enlace \ce{C-O} y
se rompe el enlace \ce{C-Br}. No hay intermedio.
\end{solution}

\begin{solution}{exo:g12:curly-arrows:11}
El cloro es más electronegativo ($3.16 - 2.55 = 0.61$ frente a
$2.96 - 2.55 = 0.41$): el carbono del clorometano es el que tiene más
carácter $\delta^+$, lo que por sí solo lo haría reaccionar más deprisa.
No es así: la velocidad depende también de la facilidad con la que el
\omterm{def:g10:electron-shells:family}{halógeno} sale con el par, y el bromo, más grande, sale con más facilidad.
\end{solution}

\begin{solution}{exo:g12:curly-arrows:12}
Si el hidrógeno se une al carbono 1: \ce{CH3-CH+-CH3}, que da
2-cloropropano \ce{CH3-CHCl-CH3}. Si se une al carbono 2:
\ce{+CH2-CH2-CH3}, que da 1-cloropropano \ce{CH2Cl-CH2-CH3}.
\end{solution}

\begin{solution}{exo:g12:curly-arrows:13}
La flecha va en sentido contrario: debe salir del dador, un \omterm{def:g10:lewis-and-shape:lone-pair}{par libre} del
\ce{OH-}, y terminar en el carbono aceptor. Y el nuevo enlace se forma
entre el oxígeno y el carbono, no entre el oxígeno y el bromo: el bromo
sale, llevándose el par del enlace \ce{C-Br}.
\end{solution}

\begin{solution}{exo:g12:curly-arrows:14}
El carbono del grupo \ce{C=O} de un ácido está unido a dos \omterm{def:g7:atoms-and-molecules:atom}{átomos} de
oxígeno electronegativos: fuertemente $\delta^+$, un aceptor. El oxígeno
del \omterm{def:g11:functional-groups:alcohol}{alcohol} lleva $\delta^-$ y dos \omterm{def:g10:lewis-and-shape:lone-pair}{pares libres}: un dador. En conjunto
se pierde agua; el \ce{-OH} del ácido se sustituye por el \ce{-O-R} del
\omterm{def:g11:functional-groups:alcohol}{alcohol}: una \omterm{def:g12:curly-arrows:categories}{sustitución}.
\end{solution}

\begin{solution}{exo:g12:curly-arrows:15}
Empujado por los \omterm{def:g9:inside-the-atom:nucleus}{electrones} del \omterm{def:g10:lewis-and-shape:multiple-bond}{doble enlace}, el par del enlace
\ce{Br-Br} se desplaza hacia el bromo lejano: el bromo cercano pasa a ser
$\delta^+$, un sitio aceptor. Primera etapa: una flecha del \omterm{def:g10:lewis-and-shape:multiple-bond}{doble enlace}
al bromo cercano, y otra del enlace \ce{Br-Br} al bromo lejano, que sale
como \ce{Br-}; queda una carga positiva en el otro carbono del antiguo
\omterm{def:g10:lewis-and-shape:multiple-bond}{doble enlace}.
\end{solution}

\begin{solution}{pb:g12:curly-arrows:1}
\textbf{1.} Eteno: \ce{H2C=CH2}, un \omterm{def:g10:lewis-and-shape:multiple-bond}{doble enlace} y cuatro \ce{C-H}.
Bromuro de hidrógeno: \ce{H-Br} con tres \omterm{def:g10:lewis-and-shape:lone-pair}{pares libres} en el bromo.

\textbf{2.} El \omterm{def:g10:lewis-and-shape:multiple-bond}{doble enlace}: su segundo par es una región rica en
\omterm{def:g9:inside-the-atom:nucleus}{electrones}.

\textbf{3.} $2.96 - 2.20 = 0.76$: el \ce{H} lleva $\delta^+$.

\textbf{4.} El \omterm{def:g7:atoms-and-molecules:atom}{átomo} de hidrógeno.

\textbf{5.} Una flecha del \omterm{def:g10:lewis-and-shape:multiple-bond}{doble enlace} al hidrógeno del \ce{HBr}; otra
del enlace \ce{H-Br} al bromo.

\textbf{6.} El carbocatión \ce{CH3-CH2+} ($+1$) y el \ce{Br-} ($-1$).

\textbf{7.} Antes: $0 + 0 = 0$; después: $+1 - 1 = 0$.

\textbf{8.} Una flecha de un \omterm{def:g10:lewis-and-shape:lone-pair}{par libre} del \ce{Br-} al carbono positivo.

\textbf{9.} Seis (tres enlaces): le falta un par para el octeto, así que
acepta un par de cualquier dador cercano.

\textbf{10.} El carbocatión \ce{CH3-CH2+}.

\textbf{11.} Se forma en la primera etapa y se consume en la segunda.

\textbf{12.} Una \omterm{def:g12:curly-arrows:categories}{adición}.

\textbf{13.} Solo el grupo: se conserva la cadena de dos carbonos, y el
\omterm{def:g10:lewis-and-shape:multiple-bond}{doble enlace} pasa a ser un enlace simple que lleva \ce{H} y \ce{Br}.

\textbf{14.} \qty{28.0}{g/mol}; \qty{108.9}{g/mol}.

\textbf{15.} $10.0 / 28.0 = \qty{0.357}{mol}$.

\textbf{16.} Eteno: $0.357 - x$; \ce{HBr}: en exceso; bromoetano: $x$.
El eteno es limitante: $x_{\max} = \qty{0.357}{mol}$.

\textbf{17.} $0.357 \times 108.9 = \qty{38.9}{g}$.

\textbf{18.} $0.75 \times 38.9 = \qty{29.2}{g}$.

\textbf{19.} Unos \qty{29}{g} de bromoetano.
\end{solution}
